\section{Trabalhos Relacionados e Artigo Base}
\label{sec:trabalhos_relacionados}

A evolu\c{c}\~{a}o da segmenta\c{c}\~{a}o sem\^{a}ntica em imagens biom\'{e}dicas pode ser categorizada em duas grandes ondas metodol\'{o}gicas: as arquiteturas convolucionais estritas e os modelos contempor\^{a}neos orientados a mecanismos de aten\c{c}\~{a}o e estruturas h\'{i}bridas.

\subsection{Segmenta\c{c}\~{a}o Convolucional em Gastroenterologia}
A introdu\c{c}\~{a}o da U-Net por Ronneberger et al. \cite{ronneberger2015unet} estabeleceu o paradigma can\^{o}nico codificador-decodificador com pontes residuais de alta resolu\c{c}\~{a}o (\textit{skip connections}). Extens\~{o}es posteriores, como a ResUNet, incorporaram blocos residuais profundos para combater a dissipa\c{c}\~{a}o do gradiente em redes de m\'{u}ltiplas camadas. N\~{a}o obstante, essas redes convolucionais homog\^{e}neas operam sob a hip\'{o}tese de invari\^{a}ncia espacial local, tornando-se incapazes de contextualizar diferen\c{c}as anat\^{o}micas globais ou de modular o ganho de aten\c{c}\~{a}o sobre les\~{o}es com contrastes atenuados perante a mucosa intestinal.

\subsection{Mecanismos de Aten\c{c}\~{a}o e Redes H\'{i}bridas Recentes}
Para mitigar a rigidez dos filtros convolucionais est\'{a}ticos, Woo et al. \cite{woo2018cbam} propuseram o \textit{Convolutional Block Attention Module} (CBAM), demonstrando que decompor a aten\c{c}\~{a}o nos eixos espacial e de canal melhora substancialmente a seletividade dos mapas de caracter\'{i}sticas intermedi\'{a}rios. No dom\'{i}nio m\'{e}dico, o advento de \textit{Vision Transformers} levou ao surgimento de redes como TransUNet \cite{chen2021transunet}, que utilizam auto-aten\c{c}\~{a}o para codificar depend\^{e}ncias de longo alcance. Contudo, transformadores puros sofrem de sobrecarga computacional massiva e demandam volumes astron\^{o}micos de dados para converg\^{e}ncia est\'{a}vel.

Como refer\^{e}ncia central de fronteira recente (2023+), o trabalho \textit{LACFormer} de Nguyen et al. \cite{lacformer2023}, publicado nos anais da \textit{British Machine Vision Conference} (BMVC 2023), investiga especificamente a converg\^{e}ncia eficiente na segmenta\c{c}\~{a}o de p\'{o}lipos. O estudo demonstra que combinar a extra\c{c}\~{a}o densa de bordas de uma CNN com m\'{o}dulos de filtragem no dom\'{i}nio espectral/aten\c{c}\~{a}o atinge acur\'{a}cias superiores com fra\c{c}\~{a}o dos par\^{a}metros de modelos puramente baseados em aten\c{c}\~{a}o global.

\subsection{Posicionamento da RAPUNet-Zeta}
Inspirada nos preceitos do LACFormer \cite{lacformer2023}, a arquitetura RAPUNet-Zeta aqui proposta adota uma abordagem estritamente engenhada: preserva a efici\^{e}ncia e a capacidade de generaliza\c{c}\~{a}o de uma CNN de grande escala (ResNet50V2), combinando-a com m\'{o}dulos de aten\c{c}\~{a}o espacial operando diretamente sobre as \textit{skip connections}. Dessa forma, elimina-se o custo proibitivo de transformadores completos, enquanto se confere \{a} rede a capacidade matem\'{a}tica de discriminar e priorizar pixels lesionais em detrimento do ru\'{i}do luminal endosc\'{o}pico.
